% Generated draft: requires mathematical and source review.
\documentclass[11pt]{article}
\usepackage[margin=25mm]{geometry}
\usepackage{fontspec}
\setmainfont{Latin Modern Roman}
\usepackage{amsmath,amssymb,mathtools,unicode-math}
\setmathfont{Latin Modern Math}
\usepackage{enumitem,needspace,xurl,hyperref}
\setlength{\parindent}{0pt}
\setlength{\parskip}{5pt}
\setlength{\emergencystretch}{3em}
\newcommand{\N}{\mathbb{N}}
\newcommand{\Z}{\mathbb{Z}}
\newcommand{\Q}{\mathbb{Q}}
\newcommand{\R}{\mathbb{R}}
\newcommand{\C}{\mathbb{C}}
\newcommand{\F}{\mathbb{F}}
\newcommand{\A}{\mathbb{A}}
\newcommand{\PP}{\mathbb{P}}
\newcommand{\E}{\mathbb{E}}
\newcommand{\eps}{\varepsilon}
\newcommand{\dd}{\,\mathrm{d}}
\begin{document}
\section*{UnsolvedMath ID: 701517}
\textbf{Source problem code:} catalogue entry
\subsection{Definitions and mathematical statement}
A group \(E\) is residually finite if, for every \(x\in E\setminus\{1\}\), there exist a finite group \(F\) and a homomorphism \(\varphi:E\to F\) such that \(\varphi(x)\ne 1\). A group is locally finite if every finitely generated subgroup is finite. A group \(Q\) has finite exponent if there exists an integer \(n\ge 1\) such that \(q^n=1\) for every \(q\in Q\).

Determine whether there exist a residually finite group \(E\) and a locally finite normal subgroup \(N\trianglelefteq E\) such that \(Q=E/N\) is infinite, finitely generated, and of finite exponent. Equivalently, ask whether such \(E,N\) can satisfy
\[
|E/N|=\infty,\qquad E/N=\langle q_1,\ldots,q_d\rangle
\quad\text{for some }d\ge 1\text{ and }q_1,\ldots,q_d\in E/N,
\]
\[
\exists n\ge 1\ \forall q\in E/N,\qquad q^n=1.
\]
There is no finite generation requirement on \(E\) or \(N\).

The supplied source also states the following conditional consequence: if no such pair \((E,N)\) exists, then there exists a hyperbolic group that is not residually finite. Here a hyperbolic group means a finitely generated group whose Cayley graph, for some finite generating set, has \(\delta\)-thin geodesic triangles for some \(\delta\ge 0\). This conditional is retained as a source assertion; its justification is not included in the supplied record.
\subsection{Short English statement}
Does there exist an infinite finitely generated group of finite exponent that is a quotient of a residually finite group by a locally finite normal subgroup? The source states that a negative answer would imply the existence of a hyperbolic group that is not residually finite.
\subsection{Sources}
\begin{itemize}
\item Local UnsolvedMath registry, permanent ID \(701517\), source code \texttt{catalogue entry}, title \emph{Kourovka Notebook Problem 18.88}; supplied complete source record, upstream revision \texttt{d371ee917ac9647fbb06cd6fc529deec6f3dda82}.
\item \emph{The Kourovka Notebook (Unsolved Problems in Group Theory)}, Problem 18.88, proposed by M. Sapir in 2014; attributed in the record to the 18th issue (2014) and the 21st issue edition (2026), \texttt{arXiv:1401.0300v48}. The supplied collection description names E. I. Khukhro and V. D. Mazurov as editors.
\item The registry labels this problem \emph{open}. This is a supplied status label, not an independently verified mathematical or literature assessment. The cited edition and the conditional consequence are reproduced from the supplied record without independent source verification.
\end{itemize}
\end{document}
